\documentclass[11pt,letterpaper]{article}
\usepackage[margin=1in]{geometry}
\usepackage[T1]{fontenc}
\usepackage{lmodern,microtype,amsmath,amssymb,amsthm,mathtools}
\usepackage[hidelinks]{hyperref}
\setlength{\emergencystretch}{2em}
\numberwithin{equation}{section}
\newtheorem{theorem}{Theorem}[section]
\newtheorem{proposition}[theorem]{Proposition}
\newtheorem{lemma}[theorem]{Lemma}
\newtheorem{corollary}[theorem]{Corollary}
\title{Self-dual forms with a pole and the smooth\\four-dimensional Poincar\'e conjecture}
\author{Edward B. Baker III}
\date{Working draft, 27 November 2026}

\begin{document}
\maketitle
\begin{abstract}
Let \(M\) be a smooth homotopy four-sphere with a conformal class \(c\).
The map \(d^-A\mapsto d^+A\) is an isometry from square-integrable
anti-self-dual two-forms to self-dual two-forms. It is the
four-dimensional counterpart of the Beurling transform, through which
the Beltrami equation gives isothermal coordinates on a surface. Its
kernel with source at a marked point \(p\) generates a three-dimensional
space of closed self-dual two-forms on \(M\setminus\{p\}\). Each possible
leading coefficient of the quartic pole occurs exactly once. If a member
has no zero, modifying its distant end and applying Gromov's relative
splitting theorem identifies \(M\) smoothly with \(S^4\). Both project
theorems have independently replayed mathematical proofs. Whether every
smooth homotopy four-sphere admits a class, point and coefficient with a
zero-free pole form remains open and is equivalent to smooth Poincar\'e.
For one compactly supported graph family on the standard sphere, we give
a uniform positive lower bound on every normalized pole direction.
For a single twisted zero circle in a marked positive cap, we characterize
a smooth boundary collar of the degree-one current by a local energy test.

\end{abstract}

\section{The problem on a fixed smooth sphere}

Let \(M\) be a closed, connected, oriented smooth manifold homotopy
equivalent to \(S^4\). Freedman's theorem identifies its homeomorphism
type \cite[Theorem~1.6]{Freedman1982}. We ask for a diffeomorphism. In Riemannian
dimension four the Hodge star splits real two-forms into two rank-three
bundles \(\Lambda^\pm\) \cite{AHS1978}. On two-forms the star depends
only on the conformal class \(c\) of the metric, and so does the \(L^2\)
inner product. All metrics are smooth and Riemannian. We normalize the
pointwise norm of two-forms so that
\(\omega\wedge\omega=|\omega|^2\,dV\) for a self-dual \(\omega\).

For a one-form \(A\) write \(dA=d^+A+d^-A\). Stokes's theorem gives
\(\|d^+A\|_{L^2}=\|d^-A\|_{L^2}\). Since \(b_2(M)=0\), every
square-integrable anti-self-dual form is \(d^-A\) for a potential \(A\in W^{1,2}\)
that is coclosed for the chosen metric, and \(A\) is unique because
\(b_1(M)=0\). Hence
\begin{equation}\label{eq:transfer}
 T_c\colon L^2(\Lambda^-)\longrightarrow L^2(\Lambda^+),\qquad
 T_c(d^-A)=d^+A,
\end{equation}
is an isometry determined by \(c\). For exact abelian curvature \(F=dA\) in real normalization it
recovers the self-dual half from the anti-self-dual half. On a surface the corresponding operator
\(f_{\bar z}\mapsto f_z\) is the Beurling transform, and the isothermal
coordinate of a conformal structure with Beltrami coefficient \(\mu\)
is obtained by inverting \(I-\mu T\). For the four-dimensional operator
compare the analysis of quasiconformal four-manifolds and of
quasiregular maps \cite{DonaldsonSullivan1989,IwaniecMartin1993}.

Fix \(p\in M\) and put \(X=M\setminus p\). For a metric \(g\in c\) let
\begin{equation}\label{eq:pole-space}
 \mathcal H_p(c)=\bigl\{\omega\in\Omega^2(X):\ d\omega=0,\
 {*}\omega=\omega,\ |\omega|_g=O\bigl(d_g(p,\cdot)^{-4}\bigr)\bigr\}.
\end{equation}
The space does not depend on the choice of \(g\) in \(c\). We call its
elements the \emph{pole forms} of \((M,c,p)\). On the round sphere, stereographic
projection from \(p\) identifies \(X\) with \(\mathbb R^4\), and the
constant self-dual forms of \(\mathbb R^4\) form a three-dimensional
space of pole forms. Each nonzero member vanishes nowhere. They play the part of the differential \(dz\) on the
Riemann sphere, which has a double pole at infinity and no zeros, and
whose primitive is the coordinate.

\begin{proposition}[A1, replayed project theorem]\label{prop:pole}
For every smooth conformal class \(c\) on \(M\) and every \(p\in M\),
the space \(\mathcal H_p(c)\) has dimension three. In oriented orthonormal normal coordinates
\(y\) at \(p\) of a metric in \(c\), each of its elements has the form
\[
 \omega=|y|^{-4}\,R_y^*\eta+O_\infty(|y|^{-2}),
\]
where \(R_y\) is the reflection in the hyperplane orthogonal to \(y\)
and \(\eta\) is a constant anti-self-dual form, and
\(\eta\mapsto\omega\) is a linear isomorphism of \(\Lambda^-_p\) onto
\(\mathcal H_p(c)\).
\end{proposition}

Here \(O_\infty(|y|^{-2})\) means that each Cartesian derivative of
order \(j\) is \(O(|y|^{-2-j})\). With \(I(y)=y/|y|^2\), the leading
term is \(I^*\eta\). The inverted chart reverses orientation.
Write \(P_g(q,p)\eta=\omega_\eta(q)\) for the solution with precisely
this leading coefficient. Against the source volume \(dV_g(p)\) the
off-diagonal kernel of \eqref{eq:transfer} is \(K_g=(2/\pi^2)P_g\).
For a fixed solution and \(\widehat g=e^{2f}g\), the covariant source
tensor satisfies \(\eta_{\widehat g}=e^{4f(p)}\eta_g\).
Thus \(\eta_g/dV_g(p)\) is independent of the representative.
Reversing orientation interchanges the points and takes the fibre adjoint,
\(K_g^{-o}(p,q)=K_g^{o}(q,p)^*\).

\begin{proof}[Proof of A1, in outline]
Set \(x=I(y)\), \(r=|x|\), and extend the end representative
\(h=|y|^{-4}g\) positively on \(X\). Its Cartesian coefficients in
\(x\) are Euclidean up to \(O_\infty(r^{-2})\).
For \(D=(d_h^*,\sqrt2\,d_h^-)\) the only flat exceptional weights in
\([-3,0]\) are the endpoints. After passing to \(rD\) on the log-radius
cylinder, the weighted Fredholm theorem and nullspace stability of
\href{https://www.numdam.org/item/ASNSP_1985_4_12_3_409_0/}
{Lockhart and McOwen (Theorem 6.2 and Lemma 7.1)} apply.
At weight \(-3/2\) the boundary term in the kernel energy identity
vanishes. A kernel one-form is then the differential of a harmonic
function tending to a constant, hence constant by the maximum principle.
In the adjoint kernel the scalar vanishes
by the maximum principle, and the remaining closed anti-self-dual form
extends over \(p\) and vanishes since \(H^2(M)=0\).
This gives invertibility throughout \((-3,0)\).

Cut off a linear primitive of the desired constant end form and solve for
a correction at weight \(-1<\delta<0\). Convolution with the flat
first-order fundamental kernel, of size \(O(r^{-3})\), sharpens the
end bounds to \(a=O_\infty(r^{-1})\), \(da=O_\infty(r^{-2})\).
For exhaustion, any form of the stipulated growth is bounded in \(x\).
The flat Hodge--de Rham operator applied to it has size \(O_\infty(r^{-3})\)
on the end. The same kernel estimate leaves a bounded entire harmonic
remainder after cutting off and subtracting the convolution, hence a
constant. Subtraction of the constructed solution with that coefficient
leaves \(O(|y|^{-2})\) at the puncture. It extends distributionally as
a closed self-dual form, hence smoothly by elliptic regularity, and
\(H^2(M)=0\) gives uniqueness. The construction uses no curvature bound.
\end{proof}

\begin{proposition}[A2, replayed project theorem]\label{prop:recognition-pole}
If some nonzero \(\omega\in\mathcal H_p(c)\) has no zero on \(X\), then
\(M\) is diffeomorphic to \(S^4\).
\end{proposition}

\begin{proof}
A nowhere-zero self-dual form is symplectic on \(X\). In the marked end
write it as \(\beta+da\) using A1. Replace \(a\) by \(\chi_R a\), where
\(\chi_R=1\) up to \(R\), vanishes past \(2R\), and has derivative
\(O(R^{-1})\). The annular error from the nonzero constant \(\beta\) is
\(O(R^{-2})\), so the modified form stays symplectic. An orthogonal map
\(Q\) and a dilation \(s>0\) put \(\beta\) in standard form in the chart
\(z=sQx\).

Cap the compact core in this chart by the complement of a ball in a
product of two equal-area spheres. The two coordinate spheres in the cap
are symplectic, have trivial normal bundles, meet positively once, and
have the same area \(A\). Excision gives \(H_2(X;\mathbb Z)=0\);
Mayer--Vietoris along the gluing three-sphere gives the two coordinate
classes as an integral basis of the capped manifold's second homology.
Its spherical periods are therefore multiples of \(A\).
Gromov's criterion and relative clause
\cite[\S2.4.A$'_1$ and A$'_2$, p.~337]{Gromov1985PseudoHolomorphic}
extend the prescribed product splitting near the coordinate spheres.
Deleting them and pasting along the common polydisc collar gives a
diffeomorphism \(X\to\mathbb R^4\) equal to \(z=sQx\) sufficiently
far out. In the inverted target coordinates the map near \(p\) is
exactly \(y\mapsto s^{-1}Qy\). It extends smoothly with smooth inverse
across the marked point.
\end{proof}

\begin{corollary}[Reformulation]\label{cor:equivalence}
The following are equivalent.
\begin{enumerate}
\item Every smooth homotopy four-sphere is diffeomorphic to \(S^4\).
\item Every smooth homotopy four-sphere \(M\) carries a conformal class
\(c\) and a point \(p\) for which some nonzero element of
\(\mathcal H_p(c)\) has no zeros.
\end{enumerate}
\end{corollary}
\begin{proof}
On the round sphere the constant forms are pole forms without zeros,
and a diffeomorphism transports the class. The converse is
Proposition~\ref{prop:recognition-pole}.
\end{proof}

In a basis \(\omega_1,\omega_2,\omega_3\) of the pole space, the
existence question asks for one fixed \(u\in S^2\) such that
\(\sum_i u_i\omega_i(q)\ne0\) for every \(q\in X\).
The class and pole may be chosen separately on each smooth type. Zero-freeness is open under smooth variation of the
conformal class, pole and nonzero coefficient. On the standard sphere,
an explicit \(U(2)\)-invariant family has invertible evaluation at every
point away from its marked pole for all \(\varepsilon>-4\).
The chiral-energy samples in \cite{Baker2026PoleFormsProposal} are
numerical diagnostics.

There are two restrictions on near-symplectic methods. In a fixed
inverted smooth chart, a positive near-symplectic form that is exactly
constant on the end is itself a pole form of a smooth compactifying
class precisely when the constant is Euclidean self-dual for the
inherited orientation. A general positive constant allows modification
far out while fixing the zero set. Conversely, near-symplecticity of a
pole form requires rank three of its intrinsic derivative at every zero.
This condition can fail on \(S^4\). Use
\href{https://doi.org/10.4310/MRL.2006.v13.n4.a6}{Taubes's local birth
(Proposition 2.1)} to obtain a transverse zero circle with unchanged
standard end. Join compact metrics realizing the zero-free and
one-circle endpoints, fixing the end germ and coefficient.
Smooth dependence and the uniform exterior estimate keep all zeros in
one compact core. At the first parameter with a zero, a rank-three zero would
persist at earlier parameters by the implicit function theorem.
This example concerns one marked direction.

This program began from the search for coordinates in dimension four
that would do what isothermal coordinates do on a surface, with twistor
theory as the guide. Linear evaluation at \(q\) induces a sphere map
\(u\mapsto\omega_u(q)/|\omega_u(q)|\) wherever the numerator is nonzero.
It gives a full fibre chart precisely at rank three. In an orthonormal
source basis, that fibre map is conformal precisely when the Gram matrix
at \(q\) is a positive scalar multiple of the identity.

The first manuscript of the program
\cite{Baker2026ChiralMaster} works with the energy
\(\int_M|W^+|^2\,dV\) of that curvature and studies the classes below
\(8\pi^2\). The two approaches meet in one question, which is open:
whether every class \(c\) on the standard \(S^4\) satisfying
\(Y(c)>0\) and \(E_+(c)<8\pi^2\) admits some pole and some
zero-free direction.

Six companion investigations concern the proofs of
the two propositions, the zero set and its first computations, the
pseudoholomorphic curves and broken fibrations attached to it, an
independent criterion through the entropy of an embedding in
\(\mathbb R^5\), finite certificates for named homotopy spheres, and
the invariants that could detect a counterexample. Their registers
record their mathematical and formal evidence separately.

\section{A definite pole triple on a compactly supported graph}

Use the oriented stereographic chart $x=(x_0,y)$ of the standard sphere
with pole at infinity. In the orthonormal flat bases
$E_i^\pm=(dx^{0i}\pm dx^{jk})/\sqrt2$, take the graph
$\omega^-=tF\omega^+$, where
\[
 g_t=a^2dx_0^2+|dy|^2,\qquad a=(1+tF)/(1-tF).
\]
The compact representative is $4g_t/(1+|x|^2)^2$.
Define $h$ to be $3/2$ times the fourth convolution power of
$\mathbf1_{[-1/2,1/2]}$, and $\rho$ the probability density
proportional to $\exp[-1/(1-(64z)^2)]$ on $|z|<1/64$, zero elsewhere.
Set $F(x)=\sum_{j=0}^4 w_j\prod_{k=0}^3(h*\rho)(x_k-c_{j,k})$.
In order, the centres are $0,e_0/2,2e_1/3,3e_2/4,4e_3/5$ and the
weights are $(16,8,4,2,1)/31$. The forms have unit Euclidean leading
coefficient in the oriented end chart.

\begin{proposition}[Marked profile A]\label{prop:marked-profile}
For $-3/4\le t\le0$ every unit leading direction $u$ satisfies
$|\omega_{t,u}(x)|_{g_t}\ge1/100$ at every $x\in\mathbb R^4$.
At $t=-3/4$ this same conformal class has $Y>2.21$ and
$0.00590\le E_+/(4\pi^2)\le1.00731$.
\end{proposition}

\begin{proof}[Proof of the pole estimate]
The flat plus field $V_tu$ equals $u+T\psi$, with minus field
$\psi=tF(u+T\psi)$. Here $T(\xi)$ is the orthogonal symbol taking the
minus component of a closed form to its plus component. The symbol is
fixed by $\iota_\xi\omega^+=\iota_\xi\omega^-$.
Let $X$ have the probability Fourier measure of $\prod_k h(x_k)$.
For fixed $x$ let $H(X)$ be $e^{ix\cdot X}$ times the mollifier's
characteristic function times $\sum_jw_j e^{-ic_j\cdot X}$. Write
$f=\mathbb E H=F(x)$, $m=\mathbb E H^2$ and
$b^2=4403963/6150400$. Positivity of the twice-mollified tensor spline
and the translation second moment give $0\le m\le\mathbb E|H|^2\le b^2$.

Put $s=|t|$. The series of successive Fourier graph corrections
converges absolutely since $sb<1$. Write its $n$th term using independent
frequency increments $X_1,\ldots,X_n$ and the multipliers at their
partial sums. Exchanging $X_1,X_2$ in terms of order at least two replaces
$T(X_1)u$ by $(T(X_1)+T(X_2))u/2$ and fixes every later partial sum.
The squared norm of this average has expectation $1/2$ since
$\mathbb E T(X)=0$. The real part of the product of the $n$ phase
factors has second moment at most $(b^{2n}+m^n)/2$. At first order
one can subtract $f$ from the real part of $H$. With the uniform
angular bound $\mathbb E(u^{\mathsf T}T(X)u)^2\le4/5$,
Cauchy--Schwarz and the geometric sum give
\[
 |u^{\mathsf T}(V_t-I)u|\le
 s\sqrt{\tfrac45\{(b^2+m)/2-f^2\}}+
 \frac{s^2\sqrt{b^4+m^2}}{2(1-sb)}.
\]
For the angular bound it suffices that two distinct coordinate components
$X^{(0)},X^{(1)}$ of one increment satisfy
$\mathbb E[(X^{(0)}X^{(1)})^2/|X|^4]>1/80$, as follows by bounding its integral
below on sixteen intervals per absolute coordinate.
The remaining functions $f$ and $m$ admit pointwise enclosures on
spatial boxes by monotonicity of $h$ and the support of $\rho$.
At $s=3/4$ the integer certificate covers $[-2,2]^4$ and bounds the
display by $99/100$. Outside, $m=0$ suffices. The right hand side
increases with $s$ and $|\omega_{t,u}|_{g_t}=(1-tF)|V_tu|$, where
$1-tF\ge1$. The analytic proof, rounding rules and executable certificate
are in the second master's note \emph{A uniform lower bound for the pole
triple on profile A}. The independent curvature and Yamabe enclosures
are in the zero-circle investigation's E05 geometry certificate.
\end{proof}

This resolves the previously undecided row of the marked experiment.
The existence of a zero-free direction for an arbitrary positive
subcritical class remains open.

\section{The boundary of the current at a single twisted circle}

Suppose the transverse input admits the prepared Taubes tube about a single
twisted zero circle $Z$ after capping on $Y=M\#\mathbb{CP}^2$. Write
$L_\infty$ for the actual square-one line, with its Kähler neighborhood.
The determinant $3\,\mathrm{PD}[L_\infty]$ in the chamber of the capped
form $\bar\omega$ supplies a finite-energy current $T$ disjoint from
$L_\infty$ \cite[Theorem~2.2 and \S7(f)]{Taubes1999SelfDualHarmonic}.
Here the chamber and the avoidance of the line use the positive-cap
calculation in investigation~03. In Taubes's orientation the outward
normal meridian pairs to $+1$ with $T$, since the canonical bundle has
degree two there, the determinant extends across the normal ball, and
$2[T]=\mathrm{PD}(c_1(\mathcal L\otimes K))$. If $d\alpha=\bar\omega$ on $Y\setminus L_\infty$,
the area of each global component is its exterior boundary degree times
$P=\int_Z\alpha$. Positivity and total degree one imply that $P>0$ and
there is one component of weight one with area $P$.

\begin{proposition}[Local density and the marked collar]\label{prop:twisted-collar}
Use the exact flat model on the normal double cover, with
$f=(x^2+y^2)/2-z^2$, $h=z(x^2+y^2)$ and
$\bar\omega=dt\wedge df+d\phi\wedge dh$, where $\phi=\arg(x+iy)$.
At $v\in Z$ measure flat distance $r_v$ from a single lift.
For a smooth nonincreasing radial cutoff $\chi$, equal to one on $[0,1]$
and zero on $[2,\infty)$, let
\[
 D_s(v)=\frac{\int_T\chi(r_v/s)\bar\omega}{b_\chi s^3},
 \qquad b_\chi=3\int_0^\infty \chi(a)a^2\,da .
\]
Take $2s<\min(R_{\rm exact},\ell/2)$, where $R_{\rm exact}$ is
the tube radius and $\ell$ is the base circle length. The full current has an embedded annular collar at $Z$ if and only if
$D_0(v)=\lim_{s\downarrow0}D_s(v)=2/3$ for all $v$. A sufficient and,
for some $s$ and some finite net, necessary test is
\[
 \max_j D_s(v_j)+\frac{\|\chi'\|_\infty P a}{b_\chi s^4}
       <\frac23+\frac1{50},
\]
where $a$ is the covering radius of the net on $Z$.
\end{proposition}

\begin{proof}
With the positive current orientation, the signed meridian reads
$p-q=1$, with $p$ equatorial and $q$ axial tangent half-planes, counted
with multiplicity. Taubes's monotonicity, tangent classification and
regularity theorem apply in the exact lifted tube.\footnote{Taubes,
\href{https://arxiv.org/pdf/math/9901142}{\emph{Geom. Topol.} 2 (1998),
221--332}, Examples 1.3--1.5, Theorem 1.8, Propositions 2.1, 3.2, 4.1
and 7.1, and \S\S6--9.}
The weighted densities equal sharp unit-ball energies on homogeneous
tangents by the choice of $b_\chi$. Tangent convergence off the axis
suffices, since rescaled area on a fixed ball is bounded and the form
vanishes linearly on the axis. Equatorial and axial energies are $2/3$
and $4/3$. Each remaining cone has energy $\delta>1/50$.
Indeed Taubes's unit link $w=(\pm3u'/2,\sqrt{c/u}\cos\tau,
\sqrt{c/u}\sin\tau,\pm u)$ gives
$3\delta=\int\sqrt{c/u+4u^2}\,|w'|\,d\tau$.
For $c\ge1/100$ both factors exceed or equal $\sqrt{c/u}$,
$u\le1$ and the interval has length at least $2\pi$.
For $c<1/100$, $\max u\ge1/\sqrt3$ and $|u'|\le2/3$ give
$u>2/5$ on an interval of length $1/2$ centred at a maximum.
There $3|u''|/2=2u/3-c/(3u^2)>1/5$ and the first factor exceeds $4/5$.

Consequently $D_0=2/3+2q+\sum\delta$. Below the displayed threshold
every tangent has $p=1$ and no axial or cone terms. The isolated-cone part in Taubes's \S6 splitting is empty and
$p\equiv1$, so every point satisfies his local-minimum criterion for regularity.
The unique analytic half-sheets glue. The deck transformation
$(t,x,y,z)\mapsto(t+\ell,x,-y,-z)$ fixes the radial parameter
$\rho=(x^2+y^2)^{1/2}$, so the graph descends to an embedded collar.
Conversely that collar has one positive equatorial tangent everywhere.
The functions $D_s$ are nondecreasing in $s$, and at fixed $s$ their
Lipschitz constant in $v$ is at most $\|\chi'\|_\infty P/(b_\chi s^4)$.
If the pointwise limit is $2/3$, compactness makes convergence uniform
along decreasing scales. Taking a sufficiently fine net at one such
scale proves necessity. The same Lipschitz and monotonicity bounds prove
sufficiency.
\end{proof}

There is also a cap-relative choice of adapted metric for which the
resulting current has one equatorial half-sheet away from finitely many
points. Choose a smaller exact tube and make compatible perturbations in
two balls centred on the axial half-planes away from the zero circle, inside the original exact region,
outside this smaller tube and disjoint from the marked cap. Taubes's annulus exclusion and
unique continuation rule out every actual regular axial germ.\footnote{Taubes,
\href{https://intlpress.com/site/pub/files/_fulltext/journals/ajm/1999/0003/0001/AJM-1999-0003-0001-a011.pdf}{\emph{Asian J. Math.} 3 (1999),
275--324}, (1.4), Proposition 2.11 and the localized Lemma 6.1.}
The form stays fixed and the chamber argument supplies a current for the
perturbed metric. Exceptional axial tangents and isolated cones still
enter Proposition~\ref{prop:twisted-collar}.

For example, add to the fixed equatorial ray along positive $x$ the
nappe $t=t_0$, $x^2+y^2=2z^2$, $z>0$, together with its deck image.
The descended local model has one regular sheet, total meridional degree
one and a cone at a single point of $Z$. Its restriction to the punctured
tube has two components. Its realization inside a globally irreducible
marked current is open. Upper bounds for the required sample integrals
on a general marked current, interior disk topology and a deformation
removing the circle relative to the cap are further questions.

\section*{Provenance and assistance}
Edward B. Baker III developed this research with assistance from OpenAI
and Anthropic language models. This draft was written for the author on
27 November 2026 by Claude Fable 5.2, at reasoning effort max as its
session was configured, from its note of the same date
\cite{Baker2026PoleFormsProposal} on the repository at
\texttt{b607bd1}. The second master integrated the mathematical proofs and
non-author readings of investigation 01 on 27 November 2026, on the
working tree based on \texttt{6a422ce}, with OpenAI Codex assistance.
The exact integration-lead model variant and effort are unavailable.
Propositions~\ref{prop:pole} and~\ref{prop:recognition-pole} are replayed
project theorems and await Lean proofs. The original relative statement
of Gromov was inspected directly. Direct inspection of Theorem~9.4.2 in
McDuff and Salamon \cite{McDuffSalamon2012JHolomorphic} remains outstanding.

The marked profile estimate was derived in the second master on the same
working tree with bounded independent OpenAI Codex calculations and
integer certificate checks. It has a mathematical replay and awaits a
formal proof. The marked collar criterion and local obstruction received
independent mathematical replays in the same investigation.

\bibliographystyle{unsrturl}
\bibliography{references}
\end{document}
